\chapter{Profile data architecture}
\label{ch:architettura_dati_profili}

The demo profiles (\texttt{small}, \texttt{medium}, \texttt{big},
\texttt{huge}, \texttt{superhuge}, \texttt{mega}) are the everyday
benchmark for the solver: the idea is to subject the engine to six
progressive school sizes (10\,$\to$\,100 classes) already loaded
with a realistic mix of constraints, so that whoever imports the
profile from the dashboard finds in front of them a \emph{stress}
scenario that exercises every constraint family the solver
handles. For years the profile was a slim pair of pickles ---
\texttt{school\_<profile>.pkl} with the rosters and
\texttt{profs\_<profile>.pkl} with the teaching assignments ---
and the import generated rooms, students and working hours on the
fly via mock helpers. The model worked as long as constraints were
few, but as constraint tables crossed the dozen mark it grew
heavy: every import had to rebuild from scratch everything the
pickle could not express, and on top of that it wiped any fixture
the user had left behind in the constraint tables. From this
chapter on, the profile source of truth is one SQLite per
profile, built once by \texttt{engine/scripts/build\_profile\_db.py}
with a deterministic seed and then imported \emph{verbatim} by the
backend.

\section{Anatomy of a profile SQLite}

Each \texttt{engine/scripts/data/<profile>/<profile>.sqlite} file
holds the same schema as \texttt{webui/data/timetable.db}, built
through \texttt{Base.metadata.create\_all} on the same
\texttt{models.py} the runtime uses. Importing the SQLite means
copying tables row-by-row into the live DB, with the projection
limited to the column intersection so the import survives partial
schema drift in either direction. To get a feel for the file's
content, open it with \texttt{sqlite3} and count rows by table
family: rosters, constraints, sites, calendar, solution.

\begin{figure}[h]
\centering
\begin{tikzpicture}[node distance=0.55cm and 0.7cm,
  every node/.style={draw, rounded corners=4pt, very thick,
                     minimum width=4.6cm, minimum height=0.95cm,
                     align=center, font=\small}]
  \node[fill=cBox!30] (an) {Rosters\\
    \textit{subjects, teachers,}\\
    \textit{school\_classes, assignments}};
  \node[fill=cBox!30, right=of an] (rm) {Rooms + sites\\
    \textit{classrooms, plessi,}\\
    \textit{plesso\_*\_rules/policies}};
  \node[fill=cBox!30, below=of an] (cs) {14 constraint tables\\
    \textit{teacher/class/room\_unavail,}\\
    \textit{coteach\_groups, gen.\,DSL}};
  \node[fill=cBox!30, below=of rm] (oh) {Weekly calendar\\
    \textit{working\_days,}\\
    \textit{working\_hour\_slots}};
  \node[fill=cBox!30, below=2cm of cs.south] (sl) {Solution (optional)\\
    \textit{solutions, lessons}};
  \draw[-{Stealth[length=2.5mm]}, thick] (an) -- (cs);
  \draw[-{Stealth[length=2.5mm]}, thick] (rm) -- (oh);
  \draw[-{Stealth[length=2.5mm]}, thick] (an) -- (rm);
  \draw[-{Stealth[length=2.5mm]}, thick] (cs) -- (sl);
  \draw[-{Stealth[length=2.5mm]}, thick] (oh) -- (sl);
\end{tikzpicture}
\caption{Anatomy of \texttt{<profile>.sqlite}: five thematic
   blocks, all populated by \texttt{build\_profile\_db.py} with a
   deterministic seed.}
\label{fig:profilo_sqlite}
\end{figure}

The file is self-sufficient: every constraint that previous
versions used to synthesise at runtime --- a
\texttt{teacher\_unavailability} marking a HARD-blocked day, a
\texttt{coteach\_groups} forcing three lab co-teaching hours, a
\texttt{plesso\_commuting\_rules} preventing a teacher from doing
the first 60 minutes at the satellite site and the next ones 800
metres away --- now lives directly as a row in its own table,
with HARD/SOFT state explicit and the same shape the SQLAlchemy
model expects in production. The build script also drops a
\texttt{manifest.json} next to the SQLite: it carries the seed,
the schedule label, and per-table row counts so the ``Stress per
profilo'' section of the Benchmarks chapter can regenerate from
scratch with one
\texttt{python -m engine.scripts.build\_profile\_db --all}.

\section{The 14 constraint tables}

The number ``fourteen'' is less daunting than the bare list
suggests, because the tables fall naturally on three axes ---
\emph{who}, \emph{when}, \emph{how} --- and the profile touches
each of them with at least one row. On the \emph{who} axis, the
first triad is teacher-centred: \texttt{teacher\_unavailability}
(HARD/SOFT cell-level state), \texttt{teacher\_mandatory\_free\_days}
(HARD: 0 hours that day, classic part-time pattern) and
\texttt{teacher\_compatible\_classes} (HARD allow-list for tight
assignments). To these add the two five-state preference tables,
\texttt{teacher\_class\_preferences} and
\texttt{teacher\_curriculum\_preferences}, which the Phase A
solver reads as HARD on the extreme branches
(\texttt{enforced} / \texttt{forbidden}) and SOFT on the
intermediate ones.

The \emph{when} axis is dominated by the three non-teacher
availability tables: \texttt{class\_unavailability},
\texttt{classroom\_unavailability} and
\texttt{logical\_unavailabilities}, the latter with parametric
DNF for expressions like ``professor X is free only on Tuesday
and Thursday afternoons''. The fourth same-axis table,
\texttt{curriculum\_logical\_constraints}, steps outside pure
time-of-day to express cross-curriculum constraints --- ``Math
and Italian not on the same day for first-years'' is the stress
sample the small profile writes by construction.

The \emph{how} axis groups co-teaching structures
(\texttt{coteach\_groups}) and the two site-aware modules:
\texttt{plesso\_commuting\_rules} (rules between site pairs,
HARD or SOFT with weight) and \texttt{plesso\_entity\_policies}
(per-teacher or per-class policies of the
\texttt{single\_plesso\_per\_day} or \texttt{single\_plesso\_total}
form). The fourteenth --- \texttt{general\_constraints} --- is
technically not a ``classic'' constraint table but holds
generic-DSL expressions (see Chapter
\ref{ch:dsl_generico_per_i_vincoli}) the post-processor evaluates
on the active solution. The small profile writes five of them,
including \texttt{forall l in lessons where
l.subject==Educazione Fisica: l.classroom.kind==palestra} which
exercises the HARD ``EduF only in gym'' rule.

\section{The capacity constraint}

The most impactful change shipped alongside the SQLite profiles
is the HARD capacity constraint: for every \texttt{Lesson},
\texttt{lesson.classroom.capacity >= lesson.school\_class.n\_students}
must hold. The rule is implemented natively in
\texttt{engine/classroom\_assignment.py}, where the
\texttt{\_can\_host} eligibility helper filters every too-small
candidate before the CP-SAT model is even built. The same rule is
also exposed as DSL pragma \texttt{classroom\_capacity\_ok()} in
the \texttt{engine/dsl\_to\_cpsat.py} compiler for solvers that
build a joint slot+room decision space (the MonolithicSolver with
a known \texttt{classroom\_for\_slot}): in that context the
pragma scans the slot space, looks up the room assigned to each
\texttt{(class, day, hour)} and forces the BoolVar to zero when
capacity falls short. The pragma is a no-op + diagnostic
``classroom data empty'' in contexts that haven't yet assigned
rooms, because then it's the Phase~C solver's job.

The input data is populated for every profile by the build
script: every class draws an \texttt{n\_students} from a truncated
gaussian over $[18, 30]$ (centred on 24, $\sigma=3.5$), and one
or two classes are deliberately forced to match exactly the
capacity of the largest standard room of the profile. The result
is a \emph{tight} constraint: the class can only land in that one
room, with everything else under pressure to schedule around. The
user sees the rule's bite at a glance in the \texttt{/classes}
tab (inline-editable ``N. studenti'' column, range 5--50) and in
\texttt{/classrooms} (inline-editable ``Capienza'' column, range
5--100); the backend remains authoritative, but the lesson
modals (\texttt{AddLessonModal}, \texttt{AddEventModal}) flag a
red warning at selection time when the (class, room) pair would
violate the rule.

\section{Special rooms and kind pragma}

Alongside capacity lives the kind constraint:
\texttt{Subject.required\_kind} is the single-row way to say
``EduF only in gym'' or ``Physics only in physics lab'', without
having to enumerate all other rooms with \texttt{forbidden}
preferences. The value is a string in the same vocabulary as
\texttt{Classroom.kind} (\texttt{standard}, \texttt{palestra},
\texttt{lab\_fisica}, \texttt{lab\_chimica},
\texttt{lab\_informatica}, \texttt{lab\_linguistico},
\texttt{biblioteca}, \texttt{aula\_speciale}); \texttt{NULL}
means ``any''. The \texttt{\_can\_host} function reads the field
from the \texttt{Lesson} (projected into the dict the backend
produces in \texttt{lessons\_for\_classroom\_step}) and rejects
every room with mismatching \texttt{kind}. The same rule is
additionally expressed as a DSL pragma in the profiles'
\texttt{general\_constraints}, so the user opening the
\texttt{/general-constraints} tab can edit it directly in the UI
without touching the model.

\section{Per-profile schedule}

A profile's stress isn't just row count: it's also the
\emph{shape} the rows take in the weekly grid. Calendar
generation, then, is not copy-paste: each profile gets its own
\texttt{working\_days} + \texttt{working\_hour\_slots} layout,
designed to put a different aspect of the solver under strain.

\begin{table}[h]
\centering
\caption{Per-profile schedules (cf.
   \texttt{engine/scripts/build\_profile\_db.py},
   \texttt{working\_schedule}).}
\label{tab:profili_orari}
\small
\begin{tabular}{lp{8.5cm}}
\toprule
\textbf{Profile} & \textbf{Weekly schedule} \\
\midrule
\texttt{small}     & 8:00--14:00 Mon--Sat, six 60\,min slots. \\
\texttt{medium}    & 8:00--14:00 Mon--Fri, 8:00--12:00 Sat. \\
\texttt{big}       & 8:00--14:00 Mon--Fri, 8:00--13:00 Sat. \\
\texttt{huge}      & 7:55--13:55 Mon--Sat, 55\,min slots with
                     a 5\,min break baked-in. \\
\texttt{superhuge} & 8:00--13:00 Mon--Fri + 14:00--17:00 Tue/Thu
                     (split-day schedule). \\
\texttt{mega}      & 8:00--14:00 Mon--Fri + 8:00--13:00 Sat +
                     14:00--17:00 Wed (one fixed afternoon). \\
\bottomrule
\end{tabular}
\end{table}

The key trick is that every schedule preserves the legacy 1--6
day numbering and 8--13 hour numbering through the
\texttt{legacy\_day\_number} and \texttt{legacy\_hour\_number}
fields of the \texttt{working\_days} / \texttt{working\_hour\_slots}
rows, so unavailability tables that reference those numbers
survive a future schedule rebuild without per-row migrations.
The \texttt{huge} profile is an interesting case: its 55-minute
slots keep \texttt{legacy\_hour\_number} 8--13, and the 5-minute
break between consecutive lessons isn't modelled as a separate
slot but as the implicit margin between \texttt{end\_time} of
slot $i$ and \texttt{start\_time} of slot $i+1$. Solvers keep
reasoning in slot units; the calendar view shows the break
naturally.

\section[Try it now]{Try it now}

To rebuild a single profile's SQLite, two shell lines suffice ---
the first invokes the script with the profile name, the second
opens the file with \texttt{sqlite3} for a sanity check:

\begin{lstlisting}[style=shell]
python -m engine.scripts.build_profile_db small
sqlite3 engine/scripts/data/small/small.sqlite \
    "select count(*) from general_constraints;"
\end{lstlisting}

\noindent The first command produces
\texttt{engine/scripts/data/small/small.sqlite} alongside the
\texttt{manifest.json} and \texttt{PICKLE\_DEPRECATED.md}; the
second prints the row count of the \texttt{general\_constraints}
table --- five, in every profile. To rebuild all six in series
(in separate subprocesses because each profile lives in its own
SQLite file and SQLAlchemy is not happy to rebind the engine
mid-process) replace the first command with
\texttt{python -m engine.scripts.build\_profile\_db --all};
the script appends each profile's \texttt{manifest.json} into
\texttt{engine/scripts/data/profiles\_manifest.json}, which the
``Stress per profilo'' section of the Benchmarks chapter reads
to populate the summary table.

\section{Compatibility with the pickle pipeline}

The \texttt{school\_<profile>.pkl} and \texttt{profs\_<profile>.pkl}
files stay in the repository for historic uses --- audit, legacy
import debugging, last-quarter run comparisons --- but they're no
longer the default for \texttt{import\_engine\_profile}: the
routine first looks for \texttt{<profile>.sqlite} in the same
folder and only uses the pickle if the SQLite is missing. A
\texttt{PICKLE\_DEPRECATED.md} note alongside each pickle
enumerates the current tables and counts for that profile, so
whoever opens the folder doesn't wonder why the import behaves
differently from what \texttt{git log} suggests: the reason is
that the SQLite is now authoritative, and it's the only file the
backend expects to find. The door to the old way stays open, but
the default path has changed.
